%% kapa: exam=1 round=29 year=2018 period=1 form=A subject=economics subject_ko=경제학원론 seq=3 number=43
%% source: Q-Net 감정평가사 기출문제 (cst00309) artlSeq=5202808 file=extracted/1008/5202808/2018년 제29회 감정평가사 1차시험 기출문제/2018년 제29회 감정평가사 1차시험 기출문제/1교시 A형.pdf
%% conversion: kordoc 4.24.51 -> tools/gen1.py (auto); review=proofread-sonnet; flags=none
\begin{question}{43}
밑줄 친 변화에 따라 2018년 Y재 시장에서 예상되는 현상으로 옳지 않은 것은? (단, 수요곡선은 우하향, 공급곡선은 우상향하며, 다른 조건은 일정하다.)
\begin{framedtext}
2017년 Y재 시장의 균형가격은 70만 원이며, 균형거래량은 500만이다. \par
2018년에 \uline{Y재 생산에 필요한 부품 가격이 상승}하였다. \par
\end{framedtext}
\begin{choices}
  \item 공급곡선은 왼쪽으로 이동한다.
  \item 균형가격은 낮아진다.
  \item 균형거래량은 줄어든다.
  \item 소비자잉여는 감소한다.
  \item 사회적 후생은 감소한다.
\end{choices}
\end{question}
\answer{2}
